\honest{Fake Utility Functions}{Eliezer Yudkowsky}{2007}

Some people have found the One Great Moral Principle from which all other values follow. I meet more of them than you do: people who know ``the amazingly simple utility function that is all you need to program into an artificial superintelligence.'' I quote none of them. \nb{No proposal is cited. In the post before, ``Fake Fake Utility Functions,'' which the book leaves out, Yudkowsky names ``happiness'' and ``genetic fitness'' as examples.}

Norman Maier advised groups not to propose solutions too soon. Friendly AI is an extremely tough problem, so people solve it extremely fast.

I may have helped cause this. Years ago I called the goal of an optimization process its ``supergoal,'' meaning ``super'' in the sense of ``parent.'' Some people seem to have heard ``the Superest Goal Ever.'' \nb{In \textsc{Creating Friendly AI} (2001), a supergoal is ``The root of a directional goal network,'' and the document speaks of supergoals in the plural.}

But a utility function need not be simple. Human values are many, with ``high Kolmogorov complexity'': a brain implements ``a thousand shards of desire,'' as anyone who has studied evolutionary psychology should know. A parent's love for a child is not derived from a child's love for a parent, or from anything else. This is a ``known fact.'' I cite no study. \nb{The post gives no evidence for it beyond the earlier posts on evolution and desire. Philosophers have long defended the parallel view about what is valuable, under the name value pluralism.}

Leave out one value, such as the value of running your own life, and you could get a ``hyperexistential catastrophe,'' a fate worse than death. \nb{The link goes to a paper by Nick Bostrom in which the word does not occur.} My example is a science-fiction story, Jack Williamson's ``With Folded Hands,'' and I link the word ``book'' to my own post warning against generalizing from fictional evidence.

The proposers do not look for objections; ``They don't know the problem is difficult.'' I know this about people I have not named. Press them on a mother's love, and they say a superintelligence that wants complexity will see how complex that love is and encourage it.

That is motivated stopping. A real maximizer of complexity would look for something more complex. It is also a fake morality. If you defend complexity by pointing out that it increases parental love, ``it means that what you really value is the parental love.''

To show that one idea is all of morality, you would have to show that everything else we think good is good only when it serves that idea. \nb{Moral philosophers who hold one basic value, Mill among them, have tried to show exactly that.} Proposers argue from good consequences instead, because that sells; one who really does reduce everything to the idea will convince no one.

In the end, only your morality reliably reproduces the decisions your morality would make. You should no more expect to compress it into a simple utility function than to compress a large file into 10 bits.
